\documentclass[11pt]{article}
\usepackage[T1]{fontenc}
\usepackage[utf8]{inputenc}
\usepackage{lmodern}
\usepackage{geometry}
\usepackage{amsmath,amssymb,amsthm,mathtools}
\usepackage{bm}
\usepackage{hyperref}
\usepackage{enumitem}
\usepackage{microtype}
\geometry{margin=1in}
\hypersetup{colorlinks=true,linkcolor=blue,citecolor=blue,urlcolor=blue}
\title{Relational Actualism III:\\Gravity, Cosmology, and the Complexity Hierarchy}
\author{Joshua F. Sandeman}
\date{Draft prepared April 2026}

\begin{document}
\maketitle

\begin{abstract}
This third paper develops the large-scale and higher-order consequences of Relational Actualism. Starting from the Local Ledger Condition and the post-actualization geometry update of the Engine of Becoming, it presents the Einstein equation as the unique macroscopic dense-limit description of the actualization graph, with departures from general relativity in sparse actualization regimes. The paper then summarizes the actualization-density account of dark-matter phenomenology, the local-$H_0$ interpretation of the Hubble tension, the daughter-spacetime picture at causal severance boundaries, and the null Bose--Marletto--Vedral prediction. It then turns to complexity: recursive coarse-graining, the Causal Firewall, life as filtered persistence, and the claim that biological and cognitive organization are higher-order motifs in the same growing causal ledger. The draft is written as the capstone of the three-paper canonical suite and incorporates the April 10--11 structural insight that the Causal Firewall is the complexity-scale analogue of topology-space exit suppression.
\end{abstract}

\section{Introduction}
Paper I introduced the kernel and Engine of Becoming. Paper II developed stable motifs, force channels, and filtered persistence. This final paper turns outward and upward: outward to geometry and cosmology, upward to the complexity hierarchy.

The central thesis is that gravitation is not a fundamental field to be quantized, but the post-actualization coarse-grained geometry of the growing ledger, while complexity is not an alien layer floating above physics, but recursive closure in the same causal graph. Sparse actualization regimes modify geometry. Deeply filtered persistence regimes generate life and agency.

\section{From Local Ledger to Einstein Geometry}
The local statement of conservation is the LLC. The macroscopic statement is the divergence-free structure of the gravitational field equations. The RA claim is that once three ingredients are in place---local conservation, causal covariance, and the local discrete curvature action---the Einstein field equations are not optional. They are the unique consistent dense-limit field equation.

In the continuum macroscopic regime the source equation is written schematically as
\begin{equation}
G_{\mu\nu}=8\pi G\, P_{\mathrm{act}}[T_{\mu\nu}],
\end{equation}
with the vacuum contribution structurally suppressed by the actualization projector.

The important conceptual point is not merely that GR is recovered. It is that GR is recovered as a post-actualization description. Geometry responds to inscribed ledger structure, not to unrealized possibility.

\section{Dense and Sparse Regimes}
RA distinguishes sharply between dense and sparse actualization regimes.
\begin{itemize}
    \item In dense regimes, the graph is sufficiently tightly woven that the continuum Einstein description is valid.
    \item In sparse regimes, the graph ceases to behave like a uniformly 3D-embeddable medium, and departures from GR become classical large-scale effects rather than Planck-suppressed quantum-gravity corrections.
\end{itemize}

This distinction is the conceptual bridge to both galactic phenomenology and the Hubble tension. The same framework that yields Einstein geometry in the dense regime predicts where and why Einstein geometry should cease to be the whole story.

\section{Actualization-Density Gravity and Dark-Matter Phenomenology}
The central sparse-regime variable is the actualization-density field $\lambda(x)$. In the weak-field, static limit the RA-modified sourcing picture separates into an accumulated-depth term and an actualization-gradient term. Schematically,
\begin{equation}
\nabla^2\Phi\;\sim\;4\pi G\bigl(\rho_A+\rho_\lambda\bigr),
\end{equation}
with the second contribution controlled by gradients of $\ln\lambda$.

The programme-level claim is that flat galactic rotation curves, galaxy-galaxy lensing profiles, and aspects of cluster phenomenology can be treated without introducing a causally silent dark matter species. The point is not that the calculations are complete in every case. The point is that the ontology changes the source term. If geometry is sourced by actualization structure rather than inert rest-mass inventory alone, then causally silent dark matter candidates lose their rationale.

\section{The Hubble Tension as a Local Expansion-Rate Effect}
In standard cosmology the Hubble tension is framed as a discrepancy between one universal $H_0$ inferred by different methods. In RA the more natural claim is that there is no single universal local expansion rate once actualization-density environments differ substantially. The inferred expansion rate is expected to depend on line-of-sight and neighborhood actualization structure.

This reframes the tension. The discrepancy is not primarily a nuisance to be eliminated; it is a phenomenon to be explained by environment-dependent coarse-graining of the actualization ledger.

\section{Causal Severance and Daughter Spacetimes}
RA treats event horizons and severance boundaries as physically special because they partition the ledger. The deeper speculative but architecturally natural proposal is that when a region reaches causal severance, a daughter spacetime can nucleate with initial conditions determined by the severed boundary data.

In the current suite this is used as a structural cosmological alternative to landscape-style anthropic reasoning. The key role of Paper III is not to overstate that programme, but to place it where it belongs: as a cosmological elaboration that flows from the severance structure of the ledger ontology.

\section{The Null BMV Prediction Revisited}
The BMV experiment belongs naturally in the gravity paper because it is the cleanest direct test of the architecture. If geometry is a Step-5 post-actualization update, then the branch-coherent Step-2 possibility structure cannot source branch-dependent gravity in advance of actualization. A positive gravity-mediated entanglement signal would therefore force a revision of the architecture at its core.

\section{Recursive Coarse-Graining and the Complexity Hierarchy}
The same ontology that yields motifs at the particle scale yields higher-order motifs under recursive coarse-graining. Let $\Pi_n$ denote the coarse-graining map from one descriptive scale to the next. Then higher-order entities arise when subgraphs achieve sufficient causal closure to be treated as effective nodes at the next level.

The hierarchy can be sketched in four tiers:
\begin{enumerate}[label=\arabic*.]
    \item Fundamental actualization events.
    \item Dissipative structures.
    \item Autopoietic systems maintaining robust closure.
    \item Agentic systems with internal generative motifs that steer future ledger growth.
\end{enumerate}

The key conceptual shift is that emergence is not treated as an epistemic complaint about intractability. It is a structural claim about stable closure under coarse-graining.

\section{The Causal Firewall}
The Causal Firewall is the complexity-scale threshold at which a shared causal ledger exists robustly enough to support open-ended biological assembly. In the present formulation it is characterized by two necessary conditions:
\begin{enumerate}[label=F\arabic*.]
    \item Sufficient redundancy of environmental record formation.
    \item A non-relativistic regime in which the causal record is not shredded by frame-dependent thermalization effects.
\end{enumerate}

At mean-field level the threshold is encoded by a percolation-type parameter
\begin{equation}
\mu=\lambda\,\tau_d\,\ell^3,
\end{equation}
with the critical regime near $\mu\sim 1$.

The April 10--11 conceptual advance was the recognition that this is the higher-order analogue of filtered exit structure in Paper II. A living system is a motif whose internal organization blocks dominant disruption channels and forces dissipation through rarer routes. In that sense the Causal Firewall is the biology-scale version of topology-space suppression.

\section{Life, Selection, and Causal Depth}
Once complexity is treated as filtered persistence, several biological claims become natural.
\begin{itemize}
    \item Life requires an environment that preserves a shared causal ledger over the assembly timescale.
    \item Natural selection tends to favor lower effective causal assembly depth at fixed function because shallow replicators reproduce faster.
    \item Boltzmann-brain-style scenarios are not merely unlikely; they conflict with the requirement of coherent causal assembly history.
\end{itemize}

These claims are bold, but they are not ontologically disconnected from the rest of the framework. They are direct descendants of the same filtered-persistence logic already encountered for particle motifs.

\section{Agency as Graph Steering}
At the highest tier, a system not only persists but actively reshapes its adjacent possible. Such a system contains a generative internal motif sufficiently rich to bias future actualization pathways. The natural action-value structure balances epistemic gain, pragmatic utility, and thermodynamic steering cost.

The point of including this in the same paper is architectural. Agency is not a metaphysical add-on imported after the physics is finished. It is a high-order closure regime in the same causal graph.

\section{What Paper III Establishes}
This final paper completes the three-paper arc.
\begin{enumerate}[label=\arabic*.]
    \item Geometry is the post-actualization coarse-grained description of the ledger.
    \item Einstein geometry is the unique dense-limit macroscopic field equation.
    \item Sparse actualization regimes generate classical departures relevant to galactic and cosmological phenomenology.
    \item The BMV protocol directly tests the proposal/inscription/geometry separation.
    \item Complexity, life, and agency are recursive closure regimes of the same causal process.
    \item The Causal Firewall is the complexity-scale expression of filtered persistence.
\end{enumerate}

\section{Conclusion}
Relational Actualism is at its most coherent when its domains are not treated as separate patchworks. Gravity is not one theory, matter another, life another still. They are different scales and regimes of one self-writing causal process.

In the dense regime the ledger looks like spacetime. In the stable local regime it looks like matter. In the recursively filtered regime it looks like life and mind. The universe is one growing actualization structure, read at different scales.

\begin{thebibliography}{99}
\bibitem{BD2010} D. M. T. Benincasa and F. Dowker, ``The scalar curvature of a causal set,'' \emph{Physical Review Letters} \textbf{104} (2010), 181301.
\bibitem{Lovelock1971} D. Lovelock, ``The Einstein tensor and its generalizations,'' \emph{Journal of Mathematical Physics} \textbf{12} (1971), 498--501.
\bibitem{SandemanRAGC} J. F. Sandeman, ``Relational Actualism Gravity and Cosmology,'' preprint, 2026.
\bibitem{SandemanRACL} J. F. Sandeman, ``Relational Actualism: The Einstein Field Equations as the Unique Macroscopic Limit of the Actualization Graph,'' preprint, 2026.
\bibitem{SandemanRADM} J. F. Sandeman, ``Relational Actualism: Dark Matter as Actualization-Density Topology,'' preprint, 2026.
\bibitem{SandemanRAHC} J. F. Sandeman, ``Algorithmic Causal Dynamics,'' preprint, 2026.
\bibitem{SandemanRACI} J. F. Sandeman, ``Relational Actualism: Life, Complexity, and the Causal Depth of Biological Systems,'' preprint, 2026.
\bibitem{SandemanRACF} J. F. Sandeman, ``The Causal Firewall,'' preprint, 2026.
\bibitem{England2013} J. L. England, ``Statistical physics of self-replication,'' \emph{Journal of Chemical Physics} \textbf{139} (2013), 121923.
\bibitem{Zurek2009} W. H. Zurek, ``Quantum Darwinism,'' \emph{Nature Physics} \textbf{5} (2009), 181--188.
\end{thebibliography}

\end{document}
